\documentclass[conference]{IEEEtran}
\IEEEoverridecommandlockouts

\usepackage{cite}
\usepackage{amsmath}
\usepackage{graphicx}
\usepackage{booktabs}
\usepackage{textcomp}

\begin{document}

\title{IMU Gesture Recognition with a Small CNN\\for a Raspberry Pi Pico}

\author{\IEEEauthorblockN{Markkus Koddala}
\IEEEauthorblockA{Student code 253101IAPM\\
IAS0360 Machine Learning for Embedded Systems\\
Tallinn University of Technology\\
markkk@taltech.ee}}

\maketitle

\begin{abstract}
This report describes a convolutional neural network that recognises four
hand gestures from accelerometer data, with model memory as the optimisation
goal. Replacing the dense head of the lab baseline with global average
pooling reduced the model from 4556 to 700 parameters without a loss in
validation accuracy (0.95 to 0.97). With int8 quantisation the weights take
0.68~KB and the TFLite file 5.0~KB, compared with 17.8~KB and 22.1~KB for the
baseline. On 20 gestures that I recorded myself the accuracy dropped to
0.62, which shows that the small shared dataset does not cover other users
well.
\end{abstract}

\begin{IEEEkeywords}
IMU, gesture recognition, CNN, quantisation, embedded systems
\end{IEEEkeywords}

\section{Method}

\subsection{Data}
The shared dataset from Lab~2.2 has four classes: \emph{negative} (no
gesture), \emph{ring} (an O shape), \emph{slope} (an L shape) and
\emph{wing} (a W shape), with 25 recordings per class. Each recording holds
about 2.1~s of raw ICM-20948 accelerometer data at about 2~kHz. The lab
preprocessor converts the counts to g, applies a 4th-order Butterworth
low-pass filter at 20~Hz, resamples each recording to 128 samples and
normalises each axis with the mean and standard deviation of the training
set. The split is made by recording, not by window, with 15/5/5 recordings
per class for training, validation and test. One whole gesture is one input
of shape $128 \times 3 \times 1$.

For my own data I recorded 5 gestures per class, 20 in total, with the Pico
and the eval board IMU. I did not have an SD card, so I changed the lab
recorder to keep each recording in RAM and send it over USB serial. The
format stays the same (\texttt{ax,ay,az} in raw counts, 2.1~s, 4200 rows).
The free-running loop sampled at 2160~Hz, so I fixed the sample period to
500~$\mu$s to get the 4200 rows that the preprocessor expects.

The shared data was recorded with the board screen facing down, so gravity
is on $+z$. I held the board upright, which puts gravity on $-x$, and drew
the gestures in the board's x--y plane. Before preprocessing, my data is
rotated into the shared frame with
\begin{equation}
\begin{bmatrix} x_s \\ y_s \\ z_s \end{bmatrix} =
\begin{bmatrix} 0 & 0 & 1 \\ 0 & 1 & 0 \\ -1 & 0 & 0 \end{bmatrix}
\begin{bmatrix} x_o \\ y_o \\ z_o \end{bmatrix}.
\end{equation}
This rotation follows from the board position, and it also gave the best
own-data accuracy of all 24 axis-aligned rotations. The own data is used
only for testing, never for training.

\subsection{Models}
The initial model is the CNN from the lab: two Conv2D + max-pooling blocks
(8 and 16 filters), Flatten, Dense(16) and a softmax output, with dropout
0.1. It has 4556 parameters, and 85\% of them are in the Dense(16) layer
after Flatten.

My own model keeps the two convolution blocks but replaces Flatten and
Dense(16) with global average pooling (GAP). Each of the 16 feature maps is
averaged over time, and the 16 values go straight to the output layer. This
gives 700 parameters. The two other architectures for Step~6 are small
versions of known networks built for the same input. The MobileNet-style
model uses a Conv2D stem and two depthwise-separable blocks (depthwise
$4 \times 1$ convolution and $1 \times 1$ pointwise convolution, 16 and 32
channels). The SqueezeNet-style model uses a Conv2D stem and two Fire
modules (a $1 \times 1$ squeeze to 4 channels, then $1 \times 1$ and
$4 \times 1$ expand convolutions with 8 channels each). Both end with GAP.

\subsection{Optimisation goal and evaluation}
I chose memory size as the optimisation goal. The Pico W has 264~KB of SRAM
and 2~MB of flash, and the model shares them with the IMU driver, the USB
stack and the TFLite Micro tensor arena. The baseline is already accurate on
the shared data, so there is little to gain in accuracy, and a smaller model
is usually faster on a microcontroller as well.

Every configuration is trained with three random seeds, because the
validation and test sets have only 20 gestures each and one gesture is
5\% of the accuracy. All models use Adam with a learning rate of
$10^{-3}$, batch size 16, and 30 epochs (60 for the GAP models, which learn
more slowly), unless stated otherwise. After training, each model is
converted to TFLite, as float32 or as full-integer int8 with post-training
quantisation, and all accuracies are measured on the TFLite model. Memory is
given as the number of parameters, the weight size (4 bytes per parameter
for float32, 1 byte for int8) and the size of the \texttt{.tflite} file.
The inference time is the time for all own gestures divided by their
number, one gesture at a time, measured on a laptop CPU with the TFLite
interpreter.

\section{Experiments}

Table~\ref{tab:opt} shows the experiments in the order they were done. The
first rows change one training parameter of the initial model at a time:
learning rate, optimiser, batch size, activation and dropout. The next rows
make the model smaller, first by using fewer filters and then with the GAP
head. The last rows apply int8 quantisation. Table~\ref{tab:arch} compares
the three architectures. The full table, with test accuracy and the spread
over seeds, is in the code repository.

\begin{table}[t]
\caption{Optimisation experiments (mean of 3 seeds)}
\label{tab:opt}
\centering
\footnotesize
\setlength{\tabcolsep}{3pt}
\begin{tabular}{lrrrrrr}
\toprule
Method & Inf. & Params & Weights & TFLite & Val. & Own \\
       & ($\mu$s) &    & (KB)    & (KB)   & acc. & acc. \\
\midrule
Initial model                  & 8  & 4556 & 17.8 & 22.1 & 0.95 & 0.40 \\
Epochs 10, lr 0.05             & 8  & 4556 & 17.8 & 22.1 & 0.93 & 0.15 \\
lr $10^{-2}$                   & 7  & 4556 & 17.8 & 22.1 & 0.95 & 0.37 \\
SGD (mom. 0.9), lr $10^{-2}$   & 7  & 4556 & 17.8 & 22.1 & 0.95 & 0.45 \\
RMSprop                        & 8  & 4556 & 17.8 & 22.1 & 0.95 & 0.43 \\
Batch size 32                  & 8  & 4556 & 17.8 & 22.1 & 0.95 & 0.32 \\
tanh activation                & 8  & 4556 & 17.8 & 22.4 & 0.95 & 0.47 \\
Dropout 0.3                    & 7  & 4556 & 17.8 & 22.1 & 0.95 & 0.27 \\
Filters 4/8, Dense(8)          & 7  & 1192 & 4.66 & 8.9  & 0.85 & 0.27 \\
\midrule
GAP head (own)                 & 7  & 700  & 2.73 & 5.9  & 0.95 & 0.53 \\
GAP, lr $10^{-2}$              & 7  & 700  & 2.73 & 5.9  & 0.97 & 0.62 \\
GAP, filters 4/8               & 6  & 224  & 0.88 & 4.1  & 0.90 & 0.30 \\
GAP, filters 4/8, lr $10^{-2}$ & 6  & 224  & 0.88 & 4.1  & 0.97 & 0.48 \\
\midrule
Initial model, int8            & 14 & 4556 & 4.45 & 10.4 & 0.95 & 0.43 \\
GAP, lr $10^{-2}$, int8        & 14 & 700  & 0.68 & 5.0  & 0.97 & 0.62 \\
GAP, 4/8, lr $10^{-2}$, int8   & 14 & 224  & 0.22 & 4.3  & 0.97 & 0.48 \\
\bottomrule
\end{tabular}
\end{table}

\begin{table}[t]
\caption{Architecture comparison (float32, mean of 3 seeds)}
\label{tab:arch}
\centering
\footnotesize
\setlength{\tabcolsep}{3pt}
\begin{tabular}{lrrrrrr}
\toprule
Architecture & Inf. & Params & TFLite & Val. & Test & Own \\
             & ($\mu$s) &   & (KB)   & acc. & acc. & acc. \\
\midrule
Initial model (lab)  & 8 & 4556 & 22.1 & 0.95 & 0.95 & 0.40 \\
Own (GAP head)       & 7 & 700  & 5.9  & 0.95 & 1.00 & 0.53 \\
MobileNet-style      & 7  & 1044 & 9.2  & 0.95 & 1.00 & 0.53 \\
SqueezeNet-style     & 7 & 628  & 9.5  & 1.00 & 1.00 & 0.67 \\
\bottomrule
\end{tabular}
\end{table}

\section{Results}

\textbf{Training parameters.} On the shared data almost every change of
learning rate, optimiser, batch size, activation or dropout gives the same
validation accuracy of 0.95. With only 20 validation gestures this is one
error, so these settings cannot be separated on validation accuracy. The
example setting from the assignment (10 epochs, learning rate 0.05) was the
only one that was clearly worse, at 0.93 on validation and 0.15 on own data.
None of these changes affect memory, because the architecture stays the
same.

\textbf{Model size.} Fewer filters in the initial model (1192 parameters)
cost 10 points of validation accuracy and made the results unstable between
seeds (standard deviation 0.13). The GAP head removed 85\% of the
parameters and kept the accuracy, because it removes the large Dense layer
instead of making the feature extractor weaker. With a learning rate of
$10^{-2}$ the GAP model reached 0.97 on validation and 1.00 on test with
700 parameters. Going down to 224 parameters still gave 0.97, but only with
the higher learning rate, and the own-data accuracy fell from 0.62 to 0.48.

\textbf{Quantisation.} int8 post-training quantisation reduced the weights
by four times without changing the validation accuracy of any model. The
effect on the file size depends on the model size. For the initial model the
TFLite file shrank from 22.1~KB to 10.4~KB. For the 700-parameter model it
only shrank from 5.9~KB to 5.0~KB, and for the 224-parameter model the int8
file (4.3~KB) was larger than the float32 file (4.1~KB). In these small
models most of the file is the graph description and the quantisation
parameters, not the weights. On the laptop the int8 models were about
twice as slow as the float32 models. This does not carry over to the Pico:
its Cortex-M0+ has no floating-point unit, so float32 arithmetic is done in
software there, and int8 is expected to be faster.

The best model for the memory goal is the GAP model with a learning rate of
$10^{-2}$ in int8: 700 parameters, 0.68~KB of weights and a 5.0~KB TFLite
file, which is 4.4 times smaller than the initial model, with higher
validation and own-data accuracy.

\textbf{Architectures.} All three architectures reach 0.95--1.00 on the
shared data. The SqueezeNet-style model has the fewest parameters (628) and
the best own-data accuracy (0.67), but its TFLite file (9.5~KB) is larger
than that of the own model (5.9~KB), because the Fire modules add more
layers and tensors to the graph. The MobileNet-style model needs more
parameters than the other two, and all three had the same inference time
on the laptop (7~$\mu$s). For the memory goal the own
GAP model is the best trade-off.

\textbf{Own data.} The accuracy on my own recordings (0.15--0.67) is much
lower than on the shared test set (0.90--1.00). With one trained own GAP model,
all negative, slope and wing gestures were recognised, apart from one
slope, but none of the five rings: two were classified as slope and three
as wing. When I mirrored the x axis of my ring recordings, the four models
recognised 19 of the 20 ring predictions, against 1 of 20 without the
mirror. So my rings differ from the training rings in the direction of
motion, most likely because I drew them in the opposite direction. My
gestures were also faster, about 1.0--1.4~s instead of almost the full
2.1~s window.

\subsection{Deployment on the Pico}\label{sec:pico}
The chosen model (GAP, learning rate $10^{-2}$, int8) was converted with
uint8 input and output, as in Lab~2.4, and run with TFLite Micro on the
Pico~W (RP2040, Cortex-M0+ at 125~MHz). The lab firmware used a different
preprocessing than the Python code (block averaging over a 4000-sample
window, no filter and no axis rotation), so I rewrote it in C++ to follow
the Python steps exactly: rotation, the two Butterworth filters, linear
resampling to 128 samples, normalisation and quantisation. To check this, a
test program runs my 20 recordings, stored in flash, through the same code.
The quantised model inputs on the Pico differ from the PC by at most one
step on single values (float32 on the Pico against float64 on the PC), and
the Pico classified 15 of the 20 gestures correctly, the same number as the
PC for this model. Two rings were labelled slope on the Pico and wing on
the PC, both wrong.

\begin{table}[h]
\caption{Chosen model on the Pico W}
\label{tab:pico}
\centering
\footnotesize
\begin{tabular}{lr}
\toprule
Model file (int8, uint8 I/O) & 5520 B \\
Tensor arena used            & 5076 B \\
Inference time per gesture   & 14.5 ms \\
Preprocessing time per gesture & 312 ms \\
Own-data accuracy (20 gestures) & 15/20 \\
Firmware flash / RAM         & 146 KB / 61 KB \\
\bottomrule
\end{tabular}
\end{table}

The whole model needs about 10.6~KB (model and arena), 4\% of the RAM of
the Pico. Inference takes 14.5~ms, but the preprocessing takes 312~ms,
because the two filters run in software floating point over 4200 samples
per axis. For this application the preprocessing, not the network, is the
part worth optimising next, for example with fixed-point filters or a
simpler decimation filter.

\section{Issues and Discussion}

The main problem is the dataset. 25 recordings per class from what looks
like one person and one session are not enough for a model that should work
for other users. The test accuracy of 1.00 says little about this; the own
data shows it. It would be better to record the gestures from several
people, in both directions and at different speeds, and to keep one person
out as the test set. I also tried training with augmented data (random time
compression and shift, rotations up to 15\textdegree{} and noise). It did
not improve the own-data accuracy and did not fix the rings, so I did not
keep it. A direction change of the ring is a mirror image, which a rotation
cannot produce.

The validation set is too small to compare training settings: most rows of
Table~\ref{tab:opt} differ by one gesture or less. Cross-validation over all
100 recordings would give more reliable numbers.

The inference times in the tables were measured on a laptop and are only
useful to compare models with each other. The Pico results are in
Section~\ref{sec:pico}.

Finally, the model expects one complete gesture in a fixed 2.1~s window
that starts when the recording starts. A real device would need a sliding
window and a way to detect the start of a gesture.

\end{document}
